\documentclass[11pt]{article}
\usepackage[margin=1in]{geometry}
\usepackage{amsmath,amssymb,amsthm,mathtools}
\usepackage{bm}
\usepackage{hyperref}
\usepackage{enumitem}
\usepackage{booktabs}
\usepackage{array}
\usepackage{longtable}
\usepackage{multirow}
\usepackage{graphicx}
\usepackage{setspace}
\usepackage{microtype}
\usepackage[T1]{fontenc}
\usepackage[utf8]{inputenc}

\title{From the Growing Ledger to Known Physics: \\ Gravity, Matter, and Coupling Structure in Relational Actualism}
\author{Joshua [Surname]}
\date{\today}

\begin{document}
\maketitle

\begin{abstract}
Relational Actualism (RA) proposes that physical reality is fundamentally composed not of objects in a pre-given spacetime, but of irreversible actualization events whose accumulation forms a growing causal ledger. A companion theory-definition paper presents the process ontology, the seven closures by which that ontology becomes a specific implemented $d=4$ theory, and the discrete substrate on which the framework rests. The present paper has a different aim. It develops the bridge from implemented $d=4$ RA to known physics.

The paper has four principal goals. First, it presents the gravity bridge, showing how local BDG action, ledger conservation, and the continuum limit yield Einstein dynamics while keeping distinct theorems, published continuum inputs, derived bridge arguments, and interpretive synthesis. Second, it presents the matter bridge: the finite stable topology universe and its systematic reading into particle classes. Third, it develops gauge and coupling structure as emergent depth-channel statistics rather than primitive sectors. Fourth, it develops a native lifetime program in which particles are modeled as self-reproducing motifs, decay is modeled as exit geometry in topology space, and the observed hierarchy of lifetimes is reinterpreted as a hierarchy of stabilization layers.

A central claim of the paper is that the implemented $d=4$ theory now supports a significantly stronger bridge architecture than earlier versions of the corpus. The local BDG coefficients are fixed by dimensional closure; the actualization threshold fixes a candidate discrimination length; that discrimination scale fixes a candidate hadronic density scale; the hadronic density scale fixes the raw direct-exit geometry of openly destabilizable motifs; and additional stabilization layers---including strangeness cost, symmetry filtering, topology protection, and branching-volume effects---stratify the longer lifetime hierarchy. In this language, what standard physics calls selection rules are reinterpreted as consistency filters on accessible exit paths in topology space.

Throughout, the paper adopts explicit status discipline. Lean-verified, computation-verified, derived, interpretive, and open-target components are separated rather than allowed to slide rhetorically into one another. The result is not a claim that every major bridge is already theorem-level. It is a claim that several of RA's strongest bridges now have enough structure to be evaluated as a coherent research program rather than as disconnected proposals.
\end{abstract}

\section*{Front-matter status box}
\noindent\textbf{This paper does:} present the strongest bridge results from implemented Relational Actualism to known physics; separate theorem-level, computation-verified, derived, and interpretive layers; and show how gravity, matter structure, interaction hierarchy, and particle lifetimes arise from the implemented $d=4$ theory.

\noindent\textbf{This paper does not:} redefine the ontology from scratch, reargue the entire kernel architecture, or survey the complete empirical prediction program. Those tasks belong to companion papers in the canonical suite.

\noindent\textbf{Status discipline:} every major bridge claim is marked, in prose, as one of the following: \textbf{LV} (Lean-verified), \textbf{CV} (computation-verified), \textbf{DR} (derived), \textbf{IN} (interpretive), \textbf{SYN} (synthesis), or \textbf{OPEN}.

\section{Introduction: From implemented theory to known physics}
Paper I defines Relational Actualism as a process ontology of irreversible actualization and presents the closure architecture by which that ontology becomes a specific implemented $d=4$ theory. The purpose of the present paper is different. It asks what known physics follows from that implemented structure, and with what epistemic status.

That question must be handled carefully. The original RA corpus was assembled during discovery, and therefore often allowed ontology, closure, bridge argument, interpretation, and prediction to appear in the same rhetorical register. That was productive during theory formation, but it makes evaluation difficult. A reader who wants to know what is theorem-level, what is exact computation, what is derived synthesis, and what remains interpretive should not be forced to infer those distinctions indirectly. This paper is therefore organized not as a catalogue of ambitious claims, but as a bridge document. Its job is to show how the implemented $d=4$ theory reaches familiar physics, while keeping every major bridge claim status-labeled and structurally separated.

The paper has four principal bridge domains. First, it presents the gravity bridge: the route from local BDG action, ledger conservation, and continuum closure to Einstein dynamics. Second, it presents the matter bridge: the route from finite topology closure to particle-class structure. Third, it presents the gauge and coupling bridge: the route from depth/sign structure to effective interaction hierarchy. Fourth, it presents the decay and persistence bridge: the route from motif renewal, density-controlled exit geometry, and $\sigma$-filtered path structure to the observed lifetime hierarchy.

This last bridge has become substantially stronger in recent work. The decay program no longer looks like a collection of phenomenological adjustments surrounding a discrete core. It now has a discernible architecture. $d=4$ closure fixes the integer substrate; that substrate fixes the selectivity threshold structure; the threshold structure fixes a candidate hadronic density scale; that density fixes the raw direct-exit geometry of openly destabilizable motifs; and additional stabilization layers then stratify the longer lifetime hierarchy through filtered exit paths. In that language, what standard physics calls selection rules are reinterpreted as consistency filters on topology-space exits. The hierarchy of lifetimes becomes a hierarchy of stabilization.

The paper is structured accordingly. Section~\ref{sec:gravity} presents the gravity bridge. Section~\ref{sec:stable-topology} presents the stable topology universe and its finite closure. Section~\ref{sec:topology-to-particles} shows how particle classes are read from that topology universe. Section~\ref{sec:gauge} develops gauge structure and force differentiation as emergent depth-channel statistics rather than primitive sectors. Section~\ref{sec:lifetimes} then develops the motif-renewal and lifetime program in full: particles as self-reproducing motifs, lifetime as exit geometry, $d=4$ closure and the hadronic density scale, direct strong decay as raw exit geometry, selection rules as stabilization layers, and pathwise exit as realized decay geometry. Section~\ref{sec:conclusion} states what the paper establishes and what remains open.

The aim, throughout, is strong but disciplined. This paper does not assume that every successful bridge is already theorem-level. It does claim that several of RA's most distinctive bridges now have enough structure to be evaluated as a coherent program rather than as disconnected proposals. That is the standard by which the sections that follow should be read.

\section{The gravity bridge}\label{sec:gravity}
\subsection{Bridge statement}
\textbf{Main status of this section:} \textbf{mixed (LV + published continuum inputs + DR + IN)}.

The gravity bridge is the strongest currently available route from implemented RA to familiar continuum physics. Its central claim is that the local BDG action, together with event-level conservation and dense-regime continuum closure, yields Einstein dynamics in the appropriate limit. The bridge should not be described as fully Lean-certified, nor should it be described as merely analogical. It is a mixed bridge result with real formal backbone, explicit external continuum inputs, and a small but nontrivial interpretive layer.

\subsection{Route I: the Lovelock chain}
The first route proceeds through the standard uniqueness logic of local, divergence-free, second-order metric dynamics. In broad outline:
\begin{enumerate}[label=\arabic*.]
    \item local combinatorial action is fixed by the implemented $d=4$ substrate;
    \item the ledger conservation law supplies the local conservation structure;
    \item the continuum limit is constrained to a second-order local form;
    \item the relevant Lovelock uniqueness theorem then identifies Einstein--Hilbert dynamics as the unique surviving term in $d=4$.
\end{enumerate}
This route relies on external published continuum theorems and should be described honestly as such.

\subsection{Route II: the RA-native BDG chain}
The second route proceeds more directly through the discrete substrate. In this route, the d=4 closure package plays a central role:
\begin{enumerate}[label=\arabic*.]
    \item the local action coefficients are fixed by dimensional closure in $d=4$;
    \item the exact second-order cancellation condition $\sum_k c_k = 0$ ensures the d'Alembertian limit;
    \item the Benincasa--Dowker continuum result then links the local action to Einstein--Hilbert structure.
\end{enumerate}
This route is more native to the RA framework, but still uses published continuum bridge results. It should therefore be described as a mixed LV/published/DR bridge.

\subsection{Ledger conservation and Bianchi structure}
The Local Ledger Condition (LLC) supplies event-level conservation. In the continuum bridge, this is read as the discrete precursor of contracted Bianchi consistency. The identification is strong but should be status-labeled as interpretive synthesis rather than theorem-level identity. The right public statement is that the LLC and Bianchi structure occupy the same role in discrete and continuum descriptions of conservation, not that the identification is already fully machine-checked.

\subsection{What the gravity bridge establishes}
The strongest justified statement is:
\begin{quote}
In the dense-regime continuum limit, implemented $d=4$ RA has a robust bridge to Einstein dynamics, supported by a mixture of discrete closure results, event-level conservation, published continuum theorems, and a small interpretive layer.
\end{quote}
This makes gravity one of the strongest bridge domains in the framework.

\section{Stable topology and the finite local matter universe}\label{sec:stable-topology}
\subsection{Bridge statement}
\textbf{Main status of this section:} \textbf{LV + CV}.

The matter bridge begins not with particle names, but with topology closure. The strongest theorem-level statement in this domain is not ``RA derives the Standard Model'' in full detail. It is that the local stable topology universe in implemented $d=4$ RA closes into a finite set of stable classes.

\subsection{Topology closure}
The relevant closure theorem, referred to in the internal architecture as L11, states that the stable local motif universe closes into exactly five classes under exhaustive single-step extension analysis. This result is paired with confinement-length structure (L10), which fixes distinct reset lengths for the relevant motif families. Together, these results show that the discrete substrate does not generate an unrestricted zoology of stable local forms. It closes.

\subsection{Why finiteness matters}
This is a powerful result independent of any later Standard Model mapping. A theory of local motif persistence that produced arbitrarily many stable types would provide little constraint on matter structure. The five-class closure result means that the implemented d=4 dynamics is restrictive enough to support a real particle-class bridge.

\subsection{What is established here}
The strongest justified statement is:
\begin{quote}
The implemented d=4 theory supports a finite and closed stable topology universe. This is the theorem-level substrate on which later particle interpretation rests.
\end{quote}

\section{From topology classes to particle classes}\label{sec:topology-to-particles}
\subsection{Bridge statement}
\textbf{Main status of this section:} \textbf{mixed (LV/CV + IN)}.

The finite topology universe does not by itself name quarks, leptons, gauge bosons, or Higgs-like motifs. The Standard Model reading is an interpretive bridge built on top of the topology closure theorem. This distinction is essential and should remain explicit.

\subsection{What the finite universe fixes}
The theorem-level closure result fixes the number of stable topology families available to the implemented theory. It also constrains the kinds of confinement, chirality, and local extension behavior that can occur in each family. This already goes well beyond loose analogy.

\subsection{What the particle mapping adds}
The particle mapping adds a systematic reading of those topology classes into familiar particle classes. That reading is not arbitrary. It is constrained by confinement behavior, chirality structure, color multiplicity, and related features of the discrete substrate. But it remains a mapping layer rather than a theorem-level identification.

\subsection{Exhaustiveness}
The strongest careful statement here is:
\begin{quote}
The topology classification is derived. The Standard Model identification is a systematic interpretive mapping built on that classification.
\end{quote}
That sentence should function as the public-facing discipline for the whole particle bridge.

\section{Gauge structure and force differentiation}\label{sec:gauge}
\subsection{Bridge statement}
\textbf{Main status of this section:} \textbf{mixed (CV + DR + IN)}.

The gauge and coupling story in RA is strongest when presented not as a finished derivation of the Standard Model gauge group, but as a constrained bridge from depth/sign structure to effective interaction hierarchy. The implemented substrate already contains a striking pattern of stabilizing and destabilizing channels, and these channels support a systematic interaction reading.

\subsection{Sign and depth structure}
The BDG vector in $d=4$ exhibits a nontrivial alternating-sign structure:
\[
(1,-1,9,-16,8).
\]
This structure already supports a native distinction between destabilizing and stabilizing channels. The same integers underwrite both the actualization threshold structure and the depth-channel hierarchy that later appears in force differentiation.

\subsection{Force differentiation}
The force-differentiation program suggests that the relevant effective channel weights take the form
\[
W_k(\mu) \sim |c_k|\,\frac{\mu^k}{k!},
\]
so that all channels are comparable at $\mu\sim 1$ while differentiating factorially at lower densities. This should not yet be described as a theorem-level derivation of gauge sectors. It is best described as a structural argument for why interaction hierarchy emerges naturally from the implemented depth structure.

\subsection{Gauge interpretation}
The mapping from sign/depth structure to gauge sectors remains a constrained interpretation rather than a finished theorem. That should be stated without embarrassment. The framework is strongest when it separates:
\begin{itemize}
    \item exact or computation-backed channel structure,
    \item derived interaction hierarchy,
    \item interpretive gauge-group identification,
    \item and open representation-theoretic targets.
\end{itemize}

\subsection{Lead-in to the lifetime program}
The most important conceptual consequence for the next section is this: forces are best understood not as primitive sectors acting on particles from outside, but as emergent depth-channel statistics shaping renewal and exit geometry for motifs. That is the natural entry point into the lifetime program.

\section{Motif renewal, selection rules, and particle lifetimes}\label{sec:lifetimes}
\subsection{Section signpost}
\textbf{Main status of this section:} \textbf{mixed (CV + DR + SYN + OPEN)}.

This section develops the bridge from the implemented d=4 theory to the observed hierarchy of particle lifetimes. The central claim is not that decay rates require ad hoc phenomenological repair layered on top of a discrete baseline. It is that the lifetime hierarchy reflects a hierarchy of accessible exit geometry in topology space. Directly destabilizable motifs decay rapidly. More persistent motifs survive because their internal organization narrows, blocks, or lengthens the accessible exit paths. In this sense, what standard physics calls selection rules are reinterpreted here as self-consistency filters on topology-space exits.

The section proceeds in five steps. First, particles are defined as self-reproducing motifs rather than primitive objects. Second, lifetime is defined as a first-exit quantity in topology space. Third, the d=4 closure package is used to motivate a derived hadronic density scale rather than a fitted one. Fourth, the direct strong-decay sector is presented as the raw exit geometry of openly destabilizable motifs. Fifth, the observed hierarchy of longer lifetimes is interpreted as a hierarchy of additional stabilization layers: strangeness cost, symmetry filtering, topology protection, and still more restrictive anomaly structure. The result is not a perturbative correction scheme but a stabilization architecture.

\subsection{Particles as self-reproducing motifs}
In Relational Actualism, a particle is not a primitive bearer of properties inhabiting an independently given spacetime. It is a persistent causal motif: a localized pattern of actualization that reproduces its effective identity under continued graph growth. The appropriate state object is therefore not merely a topology class $T$, but a fuller motif state
\[
M=(T,N,C,\sigma),
\]
where $T$ is the topology class, $N$ is the local BDG depth profile, $C$ is confinement/reset memory, and $\sigma$ is the internal label structure relevant to accessible renewal and exit channels. Topology class alone is enough to define the coarse stable family. It is not enough to define persistence, decay, or lifetime. Those require a state space rich enough to distinguish not only what a motif is, but which renewal loops are available to it and which destabilizing pathways remain open.

\subsection{Lifetime as exit geometry in topology space}
Once particles are treated as motifs, lifetime is no longer most naturally described as a primitive decay constant attached to an object. It is the expected time for a self-reproducing motif to leave its identity class. The raw lifetime variable is therefore a step-count first-exit time,
\[
\tau_P^{(\mathrm{steps})}
=
\mathbb E\!\left[\inf\{t>0:M_t\notin [M]_P\}\right].
\]
Laboratory lifetime is a late translation obtained by multiplying by the motif's own step clock. In this language, the shortest-lived resonances are those for which graph growth provides a cheap direct exit. More persistent motifs are those for which the exit geometry is narrowed or filtered.

\subsection{d=4 closure and the raw hadronic density scale}
The direct same-force decay program becomes substantially stronger if its density scale is not fitted but derived. The upstream d=4 closure package now provides the right context. The local coefficient vector is fixed by dimensional closure; the selectivity threshold is fixed by the actualization program; only d=4 satisfies the exact second-order cancellation required for the continuum wave-operator limit; and the stable topology universe is finite. Given that upstream closure network, the natural bridge question is whether the hadronic internal density is itself fixed by the threshold structure. The candidate closure is
\[
l_{\mathrm{RA}}:=\sqrt{4\Delta S^*}, \qquad
\mu_{\mathrm{QCD}}:=e^{l_{\mathrm{RA}}}.
\]
With $\Delta S^*=0.60069$ nats, one obtains $l_{\mathrm{RA}}\approx 1.550$ and therefore $\mu_{\mathrm{QCD}}\approx 4.7119$, in extremely close agreement with the earlier fine-scan best fit reported for the direct strong-decay program. The strongest current formulation is not yet theorem-complete, because the key bridge $\ln\mu = l_{\mathrm{RA}}$ remains a structurally motivated synthesis rather than a Lean theorem. But the status has nevertheless changed: the hadronic density scale is now best understood as a candidate closure quantity, not a residual free parameter.

\subsection{Direct-exit motifs and the raw exit geometry}
The direct strong-decay sector is the cleanest place where the derived density should first appear. These are motifs whose destabilizing path is openly accessible, with little or no additional exit filtering. In the later classification, they form Type~I: direct-exit motifs with $\Sigma_{\mathrm{eff}}\sim 1$. In this regime, the first-exit picture is simple. The density-controlled BDG channel weights determine the raw probability of leaving the motif identity class, and the corresponding step-count lifetime is approximately the inverse exit probability. Once converted by the motif's internal clock, this yields the laboratory lifetime scale. Direct decays are therefore the motifs for which the graph grants a short, cheap exit path.

\subsection{Selection rules as stabilization layers}\label{subsec:stabilization-layers}
The direct same-force decay baseline is not the whole lifetime story. Many motifs persist much longer than their raw direct-exit geometry would suggest. In standard language, this appears as a hierarchy of suppressed or forbidden decay channels. In the present framework, the same phenomenon is described differently: it is a hierarchy of stabilization layers that narrow, obstruct, or lengthen the accessible exit paths of a motif in topology space.

The key state variable is the full motif state
\[
M=(T,N,C,\sigma),
\]
where $T$ is topology class, $N$ is the BDG depth profile, $C$ is confinement/reset memory, and $\sigma$ is the internal label structure relevant to admissible renewal and exit channels. The function of $\sigma$ is not ornamental. It records the internal consistency conditions that determine which destabilizing transitions are actually accessible. This is why particles with the same coarse BDG profile can nevertheless have very different lifetimes: their raw exit geometry may be similar, while their accessible exit paths are not.

In this language, a selection rule is not an externally imposed prohibition on an otherwise generic decay law. It is a structural fact about the motif: certain exits are inaccessible because the motif's own internal state does not permit them. The lifetime hierarchy is therefore a stability hierarchy, generated by progressively stronger exit-path filtering.

Recent empirical synthesis suggests a five-layer classification,
\[
\text{Type I} > \text{Type II} > \text{Type III} > \text{Type IV} > \text{Type V}
\]
in accessible exit strength, and therefore the reverse in lifetime. The corresponding empirical filter strengths $\Sigma_{\mathrm{eff}}$ are reported as approximately:
\begin{itemize}
    \item Type I: $\Sigma_{\mathrm{eff}}\sim 1$,
    \item Type II: $\Sigma_{\mathrm{eff}}\sim 0.4$,
    \item Type III: $\Sigma_{\mathrm{eff}}\sim 0.1$,
    \item Type IV: $\Sigma_{\mathrm{eff}}=0$ at single-step order,
    \item Type V: $\Sigma_{\mathrm{eff}}\sim 0.002$.
\end{itemize}

The classes can be summarized as follows.

\paragraph{Type I --- Direct-exit motifs.}
These motifs possess an openly accessible destabilizing path at single-step order. Their internal label structure does not significantly narrow the direct exit. The unfiltered first-exit model therefore already captures their raw exit geometry. This class includes the direct non-strange strong decays and serves as the baseline of the stabilization hierarchy.

\paragraph{Type II --- Strangeness-cost motifs.}
These motifs possess exit paths, but the daughter structure requires flavor or strangeness rearrangement. The result is not an absolute prohibition, but a reduced accessible exit weight. This is naturally interpreted as a stabilization layer rather than an added correction. The motif persists longer because its dominant destabilizing exits are narrowed by internal consistency requirements on the daughter sector.

\paragraph{Type III --- G-parity / multi-step filtered motifs.}
These motifs admit disruption, but not direct completion of the simplest daughter channel. The destabilized intermediate cannot terminate through the shortest obvious exit and must proceed through a longer multi-step path. In this class, the content of the selection rule is pathwise rather than merely local: the first step exists, but the shortest completion does not. The $\omega$ is the canonical example.

\paragraph{Type IV --- Topology-protected motifs.}
These motifs possess no single-step exit at all. Their protection is not merely a daughter-state veto but a fact about the local BDG topology itself. The $(2,2,0,0)$ profile is the canonical case: every single-step insertion preserves positive score, so the motif is disruption-proof at single-step order. This is the RA-native interpretation of OZI-like suppression: the motif has no cheap exit.

\paragraph{Type V --- Anomaly-filtered motifs.}
These motifs are more strongly stabilized still, with effective exit only through highly constrained mixing structure. In the current empirical classification, the $\eta'$ occupies this role. This class should be regarded as the strongest presently identified stabilization layer in the meson sector, but also as the least formally closed.

The importance of the five-type classification is not only that it sorts particles into familiar qualitative categories. It shows that the lifetime hierarchy is not an arbitrary pattern laid on top of the discrete theory. It emerges as a hierarchy of progressively stronger exit-path filtering, and those filters are empirically quantifiable. The direct-exit kernel is therefore not replaced by the filtered hierarchy; it is its raw geometric substrate. The higher types are not corrections in the continuum sense. They are additional self-consistency layers by which graph growth stabilizes motifs against cheap destruction.

\begin{table}[htbp]
\centering
\caption{Five stabilization layers in the hadron sector.}
\begin{tabular}{>{\raggedright\arraybackslash}p{1.4cm} >{$}c<{$} >{\raggedright\arraybackslash}p{4.2cm} >{\raggedright\arraybackslash}p{3.0cm} >{\raggedright\arraybackslash}p{3.2cm}}
\toprule
Type & \Sigma_{\mathrm{eff}} & Stabilization mechanism & Example motifs & Role in hierarchy \\
\midrule
I & \sim 1 & direct exit accessible & $\rho,\Delta,N^*$ & raw exit geometry \\
II & \sim 0.4 & flavor / strangeness cost & $K^*,\Sigma^*,\Lambda^*$ & narrowed exit space \\
III & \sim 0.1 & shortest daughter path blocked & $\omega,a_0$ & multi-step filtered exit \\
IV & 0 & no single-step exit exists & $\phi,f_2,K_1$ & topology protection \\
V & \sim 0.002 & anomaly / mixing-dominated & $\eta'$ & strongest identified filter \\
\bottomrule
\end{tabular}
\end{table}

Type~I validates the unfiltered first-exit kernel as the raw direct-decay geometry. Types~II--V do not replace that kernel; they express additional stabilization layers that narrow, block, or lengthen the accessible exit space.

\subsection{Pathwise exit and realized decay geometry}\label{subsec:pathwise-exit}
Once a motif's shortest direct exit is blocked, one-step filtering is no longer sufficient. The relevant object is not a local channel veto alone, but the geometry of admissible exit paths through motif state space. This is the point at which the renewal program becomes explicitly pathwise.

Let
\[
M_0 \to M_1 \to \cdots \to M_n
\]
be a path $\gamma$ in motif state space beginning at the parent motif $M_0$ and ending in an exit configuration. Then the natural RA object is a pathwise exit weight of the form
\[
P(\gamma\mid M_0,\mu)
=
\prod_{j=0}^{n-1}
\Sigma_{k_j}(M_j)\,
A_{k_j}(M_j)\,
G_{k_j}(\mu)\,
B_{k_j}(M_j\to M_{j+1}),
\]
where:
\begin{itemize}
    \item $\Sigma_{k_j}(M_j)$ is the state-dependent consistency filter,
    \item $A_{k_j}(M_j)$ is structural accessibility,
    \item $G_{k_j}(\mu)$ is the density-dependent BDG channel weight,
    \item $B_{k_j}(M_j\to M_{j+1})$ is the branching weight for the state transition.
\end{itemize}

This is the natural extension of the earlier one-step factorization. It makes explicit that the motif's internal labels constrain not only whether a first disruption occurs, but what the disrupted intermediate can do next. That is why the $\omega$ could not be correctly understood at the one-step level: the problem was not the existence of disruption, but the pathwise admissibility of its completion.

The total exit probability is then a sum over admissible exit paths:
\[
P_{\mathrm{exit}}(M_0\mid\mu)
=
\sum_{\gamma\in\Gamma_{\mathrm{exit}}(M_0)}
P(\gamma\mid M_0,\mu)\,V(\gamma),
\]
where $V(\gamma)$ is the branching-volume factor associated with the path.

This last term is crucial. The recent pathwise work shows that state-dependent $\sigma$ and daughter admissibility already improve the qualitative classification substantially: $\omega$ becomes correctly Type~III, topology-protected cases become exact at single-step order, and daughter channels are restricted by baryon number, strangeness, parity, and kinematic threshold. But a remaining quantitative gap survives, especially in multi-step filtered decays. The natural interpretation is that the missing ingredient is not another ad hoc suppression factor. It is the realized exit-space volume --- the topology-space multiplicity of admissible daughter configurations reachable through a given path. In legacy continuum language, this plays the role phase space plays. In RA-native language, it is the branching volume of realized exit geometry.

This leads to a sharper statement of the selection-rule claim:
\begin{quote}
A selection rule in RA is the narrowing of realized exit geometry by internal consistency conditions on the motif and its admissible daughters.
\end{quote}

That statement has several advantages over the older vocabulary. It explains:
\begin{itemize}
    \item why direct decays are shortest-lived,
    \item why flavor-rearranging decays are slower,
    \item why G-parity-blocked decays require longer paths,
    \item and why topology-protected motifs such as the $\phi$ possess no cheap exit at all.
\end{itemize}
It also shows where the next formal work belongs. The remaining open target is not a vague phenomenological dressing of the direct kernel. It is the derivation of $V(\gamma)$, the branching-volume factor, from topology-space combinatorics and daughter admissibility.

\[
P_{\mathrm{exit}}(M\mid\mu)
=
\sum_{\gamma\in\Gamma_{\mathrm{exit}}(M)}
\left[
\prod_{j}
\Sigma_{k_j}(M_j)\,
A_{k_j}(M_j)\,
G_{k_j}(\mu)\,
B_{k_j}(M_j\to M_{j+1})
\right]
V(\gamma).
\]

\noindent\textbf{Interpretation.}
\begin{itemize}
    \item $G_k$: raw geometric opportunity,
    \item $A_k$: structural accessibility,
    \item $\Sigma_k$: consistency filtering,
    \item $B_k$: transition branching,
    \item $V(\gamma)$: realized exit-space volume.
\end{itemize}

The pathwise formalism therefore does not replace the direct-exit baseline; it completes it. The direct kernel captures the raw exit geometry of openly destabilizable motifs. The pathwise kernel captures the realized exit geometry of motifs whose accessible exits are filtered by additional stabilization layers. In this sense, the observed hierarchy of particle lifetimes is a hierarchy of how hard it is, in topology space, for a motif to complete its own destruction.

\subsection{What Section 6 establishes and what remains open}
The strongest current claims of this section are these:
\begin{itemize}
    \item particles are best modeled as self-reproducing motifs on the actualization graph;
    \item lifetime is best modeled as exit geometry in topology space;
    \item the direct strong-decay baseline is the raw exit geometry of openly destabilizable motifs;
    \item selection rules are stabilization layers that filter accessible exit paths;
    \item and the hadronic density scale is now a serious candidate closure quantity rather than a residual fit parameter.
\end{itemize}

What remains open is equally important:
\begin{itemize}
    \item the exact proof of the logarithmic bridge $\ln\mu_{\mathrm{QCD}}=l_{\mathrm{RA}}$;
    \item the derivation of $\sigma$-filters from topology plus LLC rather than empirical inference;
    \item the full pathwise daughter-state calculus;
    \item and the branching-volume theory that completes the realized exit-space geometry.
\end{itemize}
These are not patches. They are the next stabilization layers of the bridge program.

\section{Conclusion: What this paper establishes}\label{sec:conclusion}
This paper has one central aim: to show that the implemented $d=4$ version of Relational Actualism now supports a genuine bridge program to known physics, and to do so without blurring theorem-level, computation-backed, derived, and interpretive layers together.

The strongest result is not that every bridge is already closed. It is that several of the bridges are now coherent enough to be assessed as real theory components rather than exploratory gestures.

The gravity bridge is one such component. It has a mixed but robust architecture: local action closure, event-level conservation, exact second-order structure in $d=4$, and published continuum bridge theorems all converge on Einstein dynamics in the dense limit. The topology bridge is another. The finite stable topology universe is theorem-level closure; the particle-class reading built on it is systematic and constrained, even where it remains interpretive. The gauge and coupling bridge is less complete, but it already supports a nontrivial native account of interaction hierarchy as emergent depth-channel statistics.

The newest and most distinctive bridge developed here is the lifetime program. In this domain, the framework has moved beyond a loose decay phenomenology. Particles are modeled as self-reproducing motifs. Lifetime is modeled as exit geometry in topology space. The d=4 closure package fixes the discrete substrate and threshold structure. That threshold structure supports a candidate hadronic density closure. The resulting density fixes the raw direct-exit geometry of openly destabilizable motifs. Additional stabilization layers then narrow, block, or lengthen accessible exit paths, yielding the observed lifetime hierarchy. In that language, selection rules are not external prohibitions but self-consistency filters on admissible motif destruction.

Several open targets remain. The logarithmic bridge from actualization discrimination length to hadronic density should be audited further. The $\sigma$-filter hierarchy should be derived from topology and ledger structure rather than inferred empirically. The branching-volume factor should be made native, so that the realized exit-space geometry is closed at the same level as the direct-exit kernel. And the interaction and lifetime bridges should be extended beyond the hadronic strong-decay regime to the broader inter-sector structure.

These open targets do not weaken the architecture developed here. On the contrary, they are valuable because they now sit inside a cleaner map of what is already established. The bridge program is no longer a collection of unrelated claims. It has an internal order. Gravity emerges from the discrete action and its continuum limit. Matter emerges from finite topology closure. Interaction hierarchy emerges from depth-channel structure. Persistence and decay emerge from renewal, filtered exit geometry, and stabilization layers. This is not yet the whole of known physics. But it is enough to say that implemented $d=4$ RA now reaches known physics through a recognizably coherent bridge architecture.

\begin{thebibliography}{99}

\bibitem{BenincasaDowker2010}
D.~M.~T.~Benincasa and F.~Dowker,
``The scalar curvature of a causal set,''
\emph{Phys. Rev. Lett.} \textbf{104} (2010) 181301.

\bibitem{DowkerGlaser2013}
F.~Dowker and L.~Glaser,
``Causal set d'Alembertians for various dimensions,''
\emph{Class. Quantum Grav.} \textbf{30} (2013) 195016.

\bibitem{Lovelock1971}
D.~Lovelock,
``The Einstein tensor and its generalizations,''
\emph{J. Math. Phys.} \textbf{12} (1971) 498--501.

\bibitem{Yeats2025}
D.~Yeats,
``[Title of Yeats 2025 causal-diamond moment paper],''
\emph{Class. Quantum Grav.} \textbf{42} (2025) 145003.

\bibitem{RAO14}
J.~[Surname],
``RA O14: Uniqueness of the d=4 BDG action,''
internal Lean formalization and project note.

\bibitem{RAD4U02}
J.~[Surname],
``D4U02: The selectivity ceiling is at $\mu=1$,''
internal analytic proof note and certified computation package.

\bibitem{RAL11}
J.~[Surname],
``L11: Closure of the stable topology universe,''
internal Lean formalization and project note.

\bibitem{RAGS02}
J.~[Surname],
``GS02: Gauge structure in Relational Actualism,''
internal project note.

\bibitem{RAIC30}
J.~[Surname],
``IC30 and the electromagnetic coupling,''
internal project note.

\bibitem{RAGRBridge}
J.~[Surname],
``The GR bridge in Relational Actualism,''
internal audit note.

\bibitem{RATopoAudit}
J.~[Surname],
``Particle-topology audit,''
internal audit note.

\bibitem{RAMotifRenewal}
J.~[Surname],
``Motif renewal and topology-space decay in Relational Actualism,''
internal technical note.

\bibitem{RAHadronicDensity}
J.~[Surname],
``The hadronic density closure in Relational Actualism,''
internal technical note.

\bibitem{RASelectionRules}
J.~[Surname],
``Selection rules and pathwise exit in Relational Actualism,''
internal technical note.

\end{thebibliography}

\end{document}
